\section{Arranque del Nodo}

\subsection{El mecanismo \texttt{entry\_points}}

Cuando el OS ejecuta \texttt{navigate\_individual\_pva}, no ejecuta directamente el archivo Python del usuario. Ejecuta un script wrapper\footnote{Un wrapper (envolvente) es un script generado automáticamente que simplemente importa el módulo e invoca la función indicada. Actúa como puente entre el nombre del ejecutable en el sistema de archivos y la función Python real.} que \texttt{setuptools}\footnote{\texttt{setuptools} es la biblioteca estándar de Python para empaquetar y distribuir código. ROS~2 la usa para compilar paquetes Python y generar sus ejecutables durante \texttt{colcon build}.} generó automáticamente durante \texttt{colcon build}, basado en la declaración de la Terminal~\ref{lst:entry_points}.

\begin{figure}[H]
\centering
\begin{lstlisting}[style=python, caption={Declaración del ejecutable en \texttt{setup.py}.}, label={lst:entry_points}]
entry_points={
    'console_scripts': [
        'navigate_individual_pva = '
        'control_nodes.navigate_individual_pva:main'
    ],
}
\end{lstlisting}
\end{figure}

La sintaxis \texttt{nombre = modulo:funcion} le dice a \texttt{setuptools}: cuando alguien ejecute \texttt{navigate\_individual\_pva}, llama a la función \texttt{main()} del módulo \texttt{control\_nodes.navigate\_individual\_pva}. La clave \texttt{console\_scripts}\footnote{\texttt{console\_scripts} es la categoría estándar de Python para ejecutables de línea de comandos. Todo lo declarado aquí queda disponible como comando en la terminal después de instalar el paquete.} es la categoría estándar para esto. Por ello, al agregar un nuevo ejecutable al paquete es necesario recompilar con \texttt{colcon build}: el wrapper no existe hasta que se genera.

\subsection{La función \texttt{main()}}

La Terminal~\ref{lst:main_func} muestra la estructura de la función \texttt{main()}:

\begin{figure}[H]
\centering
\begin{lstlisting}[style=python, caption={Función \texttt{main()} típica de un nodo ROS~2.}, label={lst:main_func}]
def main(args=None):
    rclpy.init(args=args)
    node = NavigateIndividual()
    rclpy.spin(node)
    node.destroy_node()
    rclpy.shutdown()
\end{lstlisting}
\end{figure}

\texttt{rclpy.init(args=args)} es la primera llamada obligatoria. Inicializa la comunicación con el middleware DDS: establece la conexión con la red de nodos, configura el dominio (\texttt{ROS\_DOMAIN\_ID}), y prepara el Parameter Server\footnote{El Parameter Server en ROS~2 no es un proceso externo como en ROS~1. Es una base de datos en memoria interna a cada nodo que almacena sus parámetros con nombre, tipo y valor. Cada nodo tiene el suyo propio.} del proceso. Sin esta llamada, ninguna funcionalidad de ROS~2 está disponible: ni suscriptores, ni publicadores, ni parámetros.

\texttt{rclpy.spin(node)} entrega el control al loop de eventos\footnote{Un loop de eventos (event loop) es un patrón de programación donde el proceso espera indefinidamente a que ocurran eventos externos (llegada de mensajes, temporizadores) y ejecuta funciones cortas (callbacks) en respuesta. El proceso nunca termina por sí solo; solo sale cuando recibe una señal explícita de terminación.} y bloquea el proceso ahí hasta que llegue una señal de terminación (Ctrl+C o \texttt{SIGINT}\footnote{\texttt{SIGINT} (Signal Interrupt) es la señal que el sistema operativo envía a un proceso cuando el usuario presiona Ctrl+C en la terminal. ROS~2 captura esta señal internamente y hace que \texttt{rclpy.ok()} retorne \texttt{False}, lo que provoca que \texttt{spin()} termine ordenadamente. Sin este manejo, el proceso se mataría abruptamente sin liberar recursos.}). El código del usuario no vuelve a ejecutarse de manera lineal a partir de este punto; solo los callbacks\footnote{Un callback es una función que no se llama directamente desde el código, sino que se registra para que el sistema la invoque automáticamente cuando ocurre un evento específico. Por ejemplo, \texttt{self.odom\_callback} se registra para que ROS~2 la llame cada vez que llegue un mensaje de odometría, sin que el programador tenga que verificar manualmente la llegada de datos.} y timers registrados en \texttt{\_\_init\_\_} vuelven a activarse.

\subsection{El constructor \texttt{\_\_init\_\_}: secuencia y orden}

El constructor del nodo es donde ocurre toda la configuración inicial. No es solo inicialización de variables: es el momento en que el nodo se anuncia en la red ROS~2, declara qué parámetros acepta, y establece todos sus canales de comunicación. El orden de las operaciones importa y no es intercambiable. La Terminal~\ref{lst:node_init_full} muestra la secuencia completa:

\begin{figure}[H]
\centering
\begin{lstlisting}[style=python, caption={Constructor completo de un nodo ROS~2 con parámetros, suscriptores y timer.}, label={lst:node_init_full}]
class NavigateIndividual(Node):
    def __init__(self):
        # 1. Registro en el grafo ROS2
        super().__init__('navigate_individual')

        # 2. Declaracion de parametros (ANTES de leerlos)
        self.declare_parameter('n_rays_used', 0)
        self.declare_parameter('vmax', 1.0)

        # 3. Lectura de parametros ya resueltos
        n_rays = self.get_parameter('n_rays_used').value
        vmax   = self.get_parameter('vmax').value

        # 4. Inicializacion de estado interno
        self.pose = None
        self.odom = None

        # 5. Suscriptores
        self.create_subscription(Odometry, 'odom',
                                 self.odom_cb, 10)
        # 6. Publicadores
        self.cmd_pub = self.create_publisher(
            Twist, 'cmd_vel', 10)

        # 7. Modulo de logica (Python puro)
        self.pva = PVA(n_rays_used=n_rays, v_max=vmax)

        # 8. Timer de control (al final, cuando todo esta listo)
        self.create_timer(0.05, self.control_loop)
\end{lstlisting}
\end{figure}

\textbf{Por qué este orden.} La llamada \texttt{super().\_\_init\_\_('navigate\_individual')} debe ser la primera: registra el nodo en el grafo ROS~2\footnote{El grafo ROS~2 es la representación en tiempo real de todos los nodos activos, sus tópicos, servicios y acciones. Cuando un nodo se registra, los demás nodos pueden descubrirlo automáticamente a través de DDS sin configuración manual. Puede visualizarse con herramientas como \texttt{rqt\_graph} o \texttt{rviz2}.} con ese nombre y habilita todas las funciones de \texttt{self} (declare\_parameter, create\_subscription, get\_logger, etc.). Antes de esa llamada, el objeto existe en Python pero no tiene ninguna capacidad de ROS~2.

Los parámetros deben declararse antes de leerse. \texttt{get\_parameter} solo funciona con parámetros previamente registrados mediante \texttt{declare\_parameter}; llamarlo antes produce un error.

El timer se crea al final, una vez que todos los suscriptores, publicadores y módulos existen. Si el timer disparara antes de que \texttt{self.pva} estuviera construido, el callback intentaría usar un atributo que todavía no existe y Python lanzaría \texttt{AttributeError}\footnote{\texttt{AttributeError} es la excepción de Python que se lanza al intentar acceder a un atributo de un objeto que no existe. Por ejemplo, \texttt{self.pva.compute()} fallaría con \texttt{AttributeError} si \texttt{self.pva} aún no fue asignado.}.

\subsection{Estado interno y el problema del primer ciclo}

En programación secuencial el código fluye de arriba hacia abajo: una instrucción se ejecuta después de la anterior. En un nodo ROS~2 no hay un flujo lineal después de \texttt{spin}: el programa se vuelve orientado a eventos\footnote{La programación orientada a eventos es un paradigma donde el flujo del programa está determinado por eventos externos (mensajes, temporizadores) en lugar de por una secuencia predefinida de instrucciones. En lugar de preguntar ``¿llegó un mensaje?'' cada ciertos milisegundos (sondeo o polling), el nodo registra callbacks y el sistema lo despierta solo cuando ocurre algo relevante (notificación o event-driven). Esto ahorra CPU porque el proceso duerme mientras no hay eventos.}. El estado del nodo vive en variables de instancia (\texttt{self.pose}, \texttt{self.odom}) que los callbacks de suscripción actualizan cuando llegan mensajes, y que el timer lee cuando dispara.

El problema es que el timer puede disparar antes de que haya llegado el primer mensaje. Por eso \texttt{self.pose = None}\footnote{\texttt{None} es el valor especial de Python que representa ausencia de dato. Aquí se usa como valor centinela (sentinel): un marcador artificial que indica que una variable aún no ha recibido un valor real. Si todavía vale \texttt{None}, es porque el suscriptor no ha recibido ningún mensaje aún.}: el valor \texttt{None} actúa como centinela que señala que aún no hay información disponible. El callback de control debe verificarlo antes de operar, como muestra la Terminal~\ref{lst:none_check}:

\begin{figure}[H]
\centering
\begin{lstlisting}[style=python, caption={Verificación de datos disponibles al inicio del loop de control.}, label={lst:none_check}]
def control_loop(self):
    if self.pose is None or self.odom is None:
        return   # todavia no llego informacion, esperar

    # a partir de aqui, self.pose y self.odom son validos
    vel = self.pva.compute(self.pose, self.odom)
    self.cmd_pub.publish(vel)
\end{lstlisting}
\end{figure}

Sin esta verificación, el nodo fallaría en los primeros ciclos de control con un \texttt{AttributeError} o \texttt{TypeError}. Es una fuente frecuente de bugs en nodos ROS~2 nuevos que parece funcionar en pruebas (porque los mensajes ya llegaron) y falla al arrancar en frío.

\subsection{Declaración y lectura de parámetros en detalle}

La Terminal~\ref{lst:params_declare} ilustra el mecanismo completo:

\begin{figure}[H]
\centering
\begin{lstlisting}[style=python, caption={Declaración y lectura de parámetros en el nodo.}, label={lst:params_declare}]
self.declare_parameter('n_rays_used', 0)
self.declare_parameter('vmax', 1.0)

n_rays = self.get_parameter('n_rays_used').value  # int 5
vmax   = self.get_parameter('vmax').value          # float 0.8
\end{lstlisting}
\end{figure}

El segundo argumento de \texttt{declare\_parameter} es el valor por defecto y también define el tipo esperado. Si se declara \texttt{('vmax', 1.0)} (float) y desde la terminal se pasa \texttt{vmax:=texto}, ROS~2 lanzará un error de tipo. El tipo del default funciona como contrato\footnote{Un contrato de tipo significa que una vez declarado un parámetro como float, el Parameter Server rechaza cualquier intento de asignarle un valor de otro tipo. Esto evita errores silenciosos donde un parámetro cambia de tipo inesperadamente.}: el Parameter Server solo acepta valores del mismo tipo.

\texttt{get\_parameter('vmax').value} accede al Parameter Server y retorna el valor en el tipo Python correspondiente: \texttt{int}, \texttt{float}, \texttt{str}, \texttt{bool} o \texttt{list}. Si desde la terminal llegó \texttt{vmax:=0.8}, el YAML lo deserializó a \texttt{float(0.8)}, y \texttt{.value} lo entrega ya convertido.

\subsection{Las tres fronteras de serialización}

El valor \texttt{5} escrito en la terminal recorre tres conversiones antes de estar disponible en Python. Cada conversión es un punto donde el tipo del dato puede cambiar o perderse, lo que explica muchos errores comunes en ROS~2:

\begin{table}[H]
  \centering
  \caption{Fronteras de serialización del argumento \texttt{n\_rays\_used:=5}.}
  \label{tab:fronteras}
  \begin{tabular}{llll}
    \toprule
    Frontera & De & A & Mecanismo \\
    \midrule
    Terminal $\to$ launch & string \texttt{'5'} & string \texttt{'5'} & \texttt{split(':=')} \\
    Launch $\to$ proceso  & Python int \texttt{5}  & YAML \texttt{5} & \texttt{--ros-args -p} \\
    Proceso $\to$ nodo    & YAML \texttt{5} & Python int \texttt{5} & \texttt{.value} \\
    \bottomrule
  \end{tabular}
\end{table}

En la primera frontera no hay conversión: todo es string. La conversión a \texttt{int} ocurre en el launch file, donde el usuario la escribe explícitamente con \texttt{int(...)}. En la tercera frontera, el YAML se deserializa automáticamente y el tipo se infiere del valor declarado como default. Esto significa que el tipo final del parámetro depende de dos decisiones separadas: el \texttt{int()} en el launch file y el default en \texttt{declare\_parameter()}. Si no coinciden, pueden ocurrir errores difíciles de depurar.
